% 3. fejezet: Riasztások csoportosítása
Az előző fejezet bemutatta, hogy a behatolásérzékelők és a \ac{SIEM} rendszerek több riasztást állítanak elő, mint amennyit az elemzők érdemben fel tudnak dolgozni. Ez a fejezet azokat a módszereket tekinti át, amelyeket a kutatók és a gyártók e probléma kezelésére javasoltak. A módszereket két kérdés mentén vizsgálom: mit nyerünk velük, és milyen feltételezésekre, adatokra vagy erőforrásokra van szükségük. Mivel a dolgozatban tervezett rendszer a meglévő megoldások hiányosságaira épül, minden módszercsaládnál külön kitérek a gyengeségekre is; ezeket a fejezet végén egy összefoglaló táblázat veti össze a saját megközelítés elemeivel.

\subsection{A riasztásfeldolgozás lépései és fogalmai}\label{sec:lepesek}
A riasztások utólagos feldolgozásának szakirodalma nem egységes szóhasználatú: ugyanazt a műveletet egyes szerzők korrelációnak, mások aggregációnak vagy összevonásnak nevezik. Valeur ezért a feldolgozást egymás után kapcsolható összetevőkre bontotta, és egy közös keretben helyezték el a korábbi megoldásokat~\cite{valeur2004}. Modelljükben a riasztásokat először egységes formára hozzák, majd összevonják az ugyanazt az eseményt leíró jelzéseket, ellenőrzik, hogy a támadás sikeres lehetett-e, összekapcsolják az egy támadáshoz tartozó lépéseket, végül értékelik a kapott csoportok fontosságát. Későbbi áttekintések a korrelációs módszereket elsősorban aszerint csoportosítják, hogy hasonlóságon, előre leírt támadási forgatókönyveken, ok-okozati előfeltételeken vagy statisztikai összefüggéseken alapulnak~\cite{sadoddin2006,salah2013,mirheidari2013,kotenko2022}.

A dolgozatban három fogalmat különítek el. \emph{Csoportosításon} azt értem, amikor az összetartozó riasztásokat egyetlen egységbe -- a szakirodalomban gyakran meta-riasztásba -- fogjuk össze; ez csökkenti az elemző elé kerülő tételek számát. \emph{Rangsoroláson} azt, amikor a riasztásokat vagy csoportokat fontosság szerint sorba rendezzük; ez nem csökkenti a tételek számát, de meghatározza, mivel kezdjen az elemző. \emph{Szűrésen} pedig azt, amikor egyes riasztásokat az elemző elől teljesen elrejtünk. A szűrés a leghatékonyabb a terhelés csökkentésében, de egyben a legkockázatosabb is, mert a tévesen eltávolított riasztás valódi támadást rejthet. Ndichu és szerzőtársai áttekintése szerint a terület módszerei ezen lépések köré szerveződnek, és a legújabb munkák negyedik lépésként a nagy nyelvi modellekkel történő támogatást is bevezetik~\cite{ndichu2026}.

\subsection{A riasztások csoportosítása és korrelációja}\label{sec:csoportositas}

\subsubsection{Hasonlóság és időbeli közelség alapján}
A legegyszerűbb és a gyakorlatban legelterjedtebb megoldás azokat a riasztásokat vonja össze, amelyek néhány kulcsmezőben -- például a forrás- vagy a cél-IP-címben, a szabályban vagy a felhasználóban -- megegyeznek, és rövid időn belül követik egymást  A módszer gyors, nem igényel tanítóadatot, és könnyen magyarázható, hiszen az elemző látja, mely mező alapján kerültek egy csoportba a riasztások. Gyengesége, hogy az eredmény erősen függ a kulcsmezők és az időablak megválasztásától: túl szűk ablakkal egy támadás több csoportra esik szét, túl széles ablakkal pedig a párhuzamos, egymástól független események egy csoportba kerülnek. A mezőkre épülő szabályok ráadásul feltételezik, hogy minden riasztás tartalmazza ugyanazokat a mezőket, ami a hálózati és a gépi érzékelők riasztásainak együttes kezelésénél nem teljesül.

Landauer ezt a korlátot olyan, tartománytól független eljárással oldotta fel, amely tetszőleges, félig strukturált riasztások között képes hasonlóságot számolni, és a csoportokat időben egymást átfedő riasztásokból építi fel~\cite{landauer2022}. Kiértékelésük szerint a módszer nagyjából 80\%-kal csökkentette az elemző elé kerülő csoportok számát, és a csoportokhoz rendelt támadástípusokat mintegy 80\%-os találati és 5\%-nál kisebb téves riasztási aránnyal ismerte fel. A megközelítés azonban továbbra is az időbeli közelségre támaszkodik, ezért erősen zajos környezetben vagy egyszerre zajló támadásoknál a csoportok összemosódhatnak.

\subsubsection{Null-feltárás csoportosítással}
Julisch abból indult ki, hogy a riasztások túlnyomó többségét néhány, tartósan fennálló ok -- például egy hibásan beállított szolgáltatás vagy egy szokatlan, de jóindulatú alkalmazás -- váltja ki~\cite{julisch2003}. Vizsgálatai szerint néhány tucat ilyen kiváltó ok felelt a riasztások több mint 90\%-áért. Ha ezeket az okokat csoportosítással feltárjuk és megszüntetjük, a jövőbeli riasztások száma tartósan csökken. A módszer ereje, hogy nem tüneti kezelést ad, hanem az okot szünteti meg; gyengesége, hogy emberi értelmezést igényel, utólagos (nem valós idejű), és a valódi támadásokat nem rangsorolja, csak a jóindulatú zajt segít eltávolítani.

\subsubsection{Többlépcsős támadások összekapcsolása}
A korreláció másik iránya a támadás lépéseit kapcsolja össze: ha egy felderítést egy sikeres bejelentkezés, majd jogosultságszerzés követ ugyanazon a gépen, ezek együtt sokkal többet jelentenek, mint külön-külön. Erre a célra a szakirodalom támadási forgatókönyveket, a támadás szakaszait követő állapotgépeket~\cite{wilkens2021}, a riasztásokból tanult támadási gráfokat~\cite{nadeem2021} és a korreláció utáni elemzést~\cite{shittu2015} egyaránt javasolt. Közös gyengeségük, hogy előre rögzített tudást -- forgatókönyveket, a lépések sorrendjét vagy tanítóadatot -- igényelnek, így a modellben nem szereplő támadási lánc rejtve maradhat, a fenntartásuk pedig jelentős szakértői munkát kíván.

\subsubsection{Gépi tanulásra és gráfokra épülő csoportosítás}
Az újabb munkák a csoportosítást tanuló modellekkel végzik. A DeepCASE az egy géphez tartozó eseménysorozatokból tanulja meg, mely korábbi események adnak magyarázatot egy újabb riasztásra, és a hasonló helyzeteket egyszerre, egy döntéssel kezelteti az elemzővel~\cite{vanede2022}. A szerzők szerint az események 86,72\%-át tudta így automatikusan kezelni, és az elemzői munkát 90,53\%-kal csökkentette. Az AlertBERT a riasztásokat egy maszkolt nyelvi modellel alakítja számvektorrá, majd sűrűségalapú csoportosítást végez; a szerzők szerint zajos környezetben és egyidejű támadások esetén megbízhatóbb, mint az időkülönbségre épülő csoportosítás~\cite{karner2026alertbert}. Eckhoff a riasztásokat közös jellemzőik -- például IP-cím vagy felhasználó -- mentén gráffá kötik össze, és a gráfrészletek hasonlóságát tanuló modellel vizsgálják~\cite{eckhoff2025}.

Ezek a módszerek pontosabbak a kézi szabályoknál, de áruk van. Tanítóadatot vagy legalább egy betanítási időszakot igényelnek, a környezet változásakor újra kell tanítani őket, és döntéseik nehezebben magyarázhatók az elemző számára. Kiértékelésük ráadásul jellemzően a csoportok tisztaságára vagy a csökkentés mértékére irányul, arra kevésbé, hogy a csoportok közül melyik a legfontosabb.

\begin{table}[H]
    \centering
    \caption{A riasztás-csoportosítási módszerek összehasonlítása (Forrás: saját összeállítás)}
    \label{tab:csoportositas}
    \small
    \begin{tabular}{|L{2.8cm}|L{2.4cm}|L{4.0cm}|L{4.0cm}|}
        \hline
        \textbf{Módszercsalád} & \textbf{Példa} & \textbf{Erősség} & \textbf{Gyengeség} \\ \hline
        Kulcsmezők és időablak & \cite{valeur2004} & Gyors, tanítás nélküli, átlátható & Az ablak és a mezők megválasztásától függ; eltérő formátumú riasztásoknál nehéz \\ \hline
        Tartomány\-független hasonlóság & \cite{landauer2022} & Tetszőleges riasztásformára működik & Az időbeli közelségre épül; zajban a csoportok összemosódnak \\ \hline
        Ok-feltárás & \cite{julisch2003} & A zaj okát szünteti meg & Utólagos, emberi értelmezést kíván, nem rangsorol \\ \hline
        Többlépcsős korreláció & \cite{wilkens2021,nadeem2021} & A támadási láncot egyben mutatja & Előre rögzített tudást vagy tanítóadatot igényel \\ \hline
        Gépi tanulás, gráf & \cite{vanede2022,karner2026alertbert,eckhoff2025} & Pontos, zajtűrő & Tanítást igényel, nehezen magyarázható, nem a fontosságot méri \\ \hline
    \end{tabular}
\end{table}

\subsection{A riasztások rangsorolása}\label{sec:rangsorolas}

\subsubsection{Rangsorolás a súlyosság szerint}
A legtöbb behatolásérzékelő és \ac{SIEM} rendszer minden szabályhoz súlyossági szintet rendel, és az elemzők ezt használják a sorrend kialakítására. A módszer előnye, hogy semmilyen többletadatot nem igényel; hátránya, hogy a súlyosságot a szabály írója a szabály megírásakor, a konkrét környezet ismerete nélkül határozza meg. Ugyanaz a szabály ugyanazt a súlyosságot kapja egy tesztgépen és egy kritikus adatbázis-kiszolgálón, és akkor is, ha a megcélzott szoftver nem is fut az adott gépen. A súlyossági szintek emellett durvák: néhány fokozatba sok ezer riasztás kerül, amelyeken belül nincs további sorrend. Uetz és szerzőtársai mérése szerint a csak súlyosság szerinti sorrend rangsor-pontossága (\ac{AUROC}) átlagosan 0,72 volt~\cite{uetz2026rba}, Ndichu és szerzőtársai áttekintése pedig kifejezetten kockázatként írja le, ha egy csapat kizárólag erre hagyatkozik~\cite{ndichu2026}.

\subsubsection{Többszempontú döntési módszerek}
A többszempontú döntéstámogató módszerek több mérőszámot vonnak össze egyetlen pontszámmá. Anuar és szerzőtársai a páros összehasonlításon alapuló \ac{AHP}-t használták arra, hogy az érintett eszköz kritikusságát, helyettesíthetőségét és más tulajdonságait súlyozzák, és ebből kockázati indexet számoljanak~\cite{anuar2012}. Kiértékelésükben ezzel az incidensek mindegyikét sorba lehetett állítani, míg a sérülékenységek szabványos pontozása (\ac{CVSS}) alapján csak 17,23\%-ukat. Alsubhi és szerzőtársai fuzzy logikával kombinálták többek között a támadás alkalmazhatóságát, az áldozat fontosságát, az érzékelő állapotát és a korábbi riasztásokkal való kapcsolatot~\cite{alsubhi2008,alsubhi2011fuzmet}.

E módszerek erőssége, hogy átláthatók és szakértői tudással közvetlenül hangolhatók. Gyengeségük, hogy a súlyokat és a szabályokat jellemzően egy szakértő adja meg, ennek hatását ritkán vizsgálják érzékenységvizsgálattal, és kiértékelésük többnyire régi, mesterséges
adathalmazokon történt. Így nem tudjuk, mennyit érnek egy korszerű, erős összehasonlítási alaphoz képest. Jalalvand és szerzőtársai áttekintése a rangsorolás szempontjait és módszereit rendszerezi, és ugyancsak azt jelzi, hogy a szempontok bősége mellett kevés az összevethető, megismételhető kiértékelés~\cite{jalalvand2024}.

\subsubsection{Környezeti információ és riasztás-ellenőrzés}
Egy riasztás jelentősége nagymértékben attól függ, hol történik. Porras és szerzőtársai már 2002-ben olyan rangsorolást javasoltak, amely a szervezet működése szempontjából fontos gépeket és szolgáltatásokat, valamint a hálózat felépítését is figyelembe veszi~\cite{porras2002}. Kruegel és szerzőtársai a riasztás-ellenőrzést vizsgálták: azt, hogy egy támadási kísérlet sikeres lehetett-e az adott célponton~\cite{kruegel2004verification}. Rámutattak, hogy a környezeti információ önmagában nem elég, mert az ellenőrzés csak akkor megbízható, ha a gépekről tárolt adatok naprakészek. Njogu és szerzőtársai a hálózat sérülékenységeit vetették össze a riasztásokkal, és azokat a riasztásokat, amelyek nem érintettek meglévő sérülékenységet, alacsonyabb prioritásúnak tekintették~\cite{njogu2012}. Hasonló irányt követnek a tudásalapú és a kontextusalapú szűrők is~\cite{meng2015,chergui2020}.

A környezetre épülő módszerek legfőbb gyengesége, hogy jellemzően \emph{szűrnek}: a nem illeszkedő riasztást eltávolítják vagy hátrasorolják. Ha a nyilvántartás hiányos -- egy gépről nem tudjuk, hogy fut rajta egy sérülékeny szolgáltatás --, a valódi támadás is kieshet. Sommestad és Franke kísérletében a hálózati információra épülő szűrés valóban jelentősen csökkentette a riasztások számát, de a valódi támadásokhoz tartozó riasztások egy részét is eltávolította~\cite{sommestad2015}. További gyengeség, hogy e munkák többsége a környezeti információt egyetlen, magában álló rangsorolóban használja, és nem méri, mennyit tesz hozzá egy olyan módszerhez, amely a riasztások folyamából már önmagában is jó sorrendet ad.

\subsubsection{Kockázatalapú riasztás-pontozás}
A kockázatalapú riasztás-pontozás (\ac{RBA}) nem az egyes riasztásokat, hanem az egy géphez vagy felhasználóhoz kapcsolódó, időben halmozódó kockázatot értékeli. Az iparban a Splunk vezette be ezt a szemléletet, ahol a riasztások kockázati pontokat adnak az érintett eszközhöz vagy személyhez, és az elemző csak akkor kap jelzést, ha az összeg egy küszöböt átlép~\cite{splunkrba}. A gyártói megvalósítás az eszközök és személyek fontosságát is beszámíthatja, a hatásáról azonban nincs nyilvános, megismételhető mérés.

Uetz és szerzőtársai a módszer nyílt, tanítóadatot nem igénylő változatát dolgozták ki és mérték fel nyolc adathalmazon~\cite{uetz2026rba}. A pontszám öt, a riasztásokból közvetlenül számolható részből áll: a szabály súlyosságából, a riasztások rövid időn belüli halmozódásából, a változatosságukból, valamint a ritkaságukból és a szokásos mintától való eltérésükből. Ezzel a rangsor-pontosság átlagosan 0,92-re nőtt a csak súlyosság szerinti 0,72-höz képest. Ez a módszer a dolgozat legerősebb összehasonlítási alapja: nyilvános, megismételhető, és nem igényel tanítást. Gyengesége ugyanakkor, hogy minden információja a riasztások folyamából származik. Nem tudja, hogy az érintett gép mennyire fontos, fut-e rajta a támadás által kihasználható szoftver, és mennyire látják a megfigyelőeszközök; egy kritikus, de ritkán riasztó kiszolgálón történő egyetlen, célzott támadás így hátrébb kerülhet egy zajos tesztgép riasztásainál.

\subsubsection{Tanuló és származáskövető rangsorolás}
A tanuló módszerek a fontosságot adatból becslik. Az AlertPro a riasztások környezeti jellemzőit megerősítéses tanulással dolgozza fel, és az elemzők visszajelzéseiből tanulja meg a többlépcsős támadások lépéseinek előresorolását~\cite{wang2024alertpro}. A NoDoze minden riasztáshoz felépíti az érintett folyamatok és fájlok ok-okozati gráfját, és azt értékeli, mennyire szokatlan ez a gráf a szervezetben korábban látottakhoz képest~\cite{hassan2019}. A szerzők 364 riasztáson végzett kiértékelése szerint a téves riasztások száma 86\%-kal csökkent. A RAPID ugyanezt a származáskövetést valós idejűvé tette, és a már folyamatban lévő vizsgálatokat is figyelembe veszi~\cite{liu2022rapid}. Hasonló, megerősítéses tanulásra vagy tanuló modellekre épülő megoldások a legfrissebb irodalomban is megjelennek~\cite{homayoun2026,guo2026}.

E módszerek gyengesége a tanítóadat iránti igény: az elemzők címkéi, visszajelzései vagy hosszú megfigyelési időszak nélkül nem használhatók, a címkék pedig a valós SOC-okban ritkák és következetlenek. A származáskövető megoldások ráadásul részletes, gépi szintű eseménynaplót igényelnek, amely a hálózati behatolásérzékelők riasztásaihoz nem áll rendelkezésre. Végül a tanult súlyok nehezen magyarázhatók, és a környezet megváltozásakor újratanítást igényelnek.

\begin{table}[H]
    \centering
    \caption{A riasztás-rangsorolási módszerek összehasonlítása (Forrás: saját összeállítás)}
    \label{tab:rangsorolas}
    \small
    \begin{tabular}{|L{2.5cm}|L{2.3cm}|L{1.6cm}|L{1.6cm}|L{4.8cm}|}
        \hline
        \textbf{Módszercsalád} & \textbf{Példa} & \textbf{Környe\-zet} & \textbf{Tanítás} & \textbf{Fő gyengeség} \\ \hline
        Súlyosság & gyártói szabályok & nem & nem & Durva, a környezettől független \\ \hline
        Többszempontú (\ac{AHP}, fuzzy) & \cite{anuar2012,alsubhi2011fuzmet} & igen & nem & Szakértői súlyok érzékenységvizsgálat nélkül; régi adathalmazokon mérve \\ \hline
        Környezet és ellenőrzés & \cite{porras2002,kruegel2004verification,njogu2012} & igen & nem & Szűr: hiányos nyilvántartásnál valódi támadást is eltávolít \\ \hline
        Kockázatalapú (\ac{RBA}) & \cite{uetz2026rba,splunkrba} & nyilvános mérésben nem & nem & Csak a riasztások folyamából dolgozik \\ \hline
        Tanuló, származáskövető & \cite{wang2024alertpro,hassan2019,liu2022rapid} & részben & igen & Címkét vagy részletes gépi naplót igényel; nehezen magyarázható \\ \hline
    \end{tabular}
\end{table}

\subsection{Agent-alapú és nagy nyelvi modellekre épülő megközelítések}\label{sec:agentek}

\subsubsection{Klasszikus többágenses rendszerek}
Az agent-alapú felépítés nem új a behatolásérzékelésben. A többágenses rendszerekben önálló, egymással üzenetekkel együttműködő programrészek végzik az adatgyűjtést, az elemzést és a döntést; ez jól illeszkedik az elosztott hálózatokhoz, és lehetővé teszi, hogy az egyes feladatokat külön, cserélhető összetevők lássák el~\cite{bougueroua2021}. A TRINETR például külön agentekkel gyűjtötte össze a riasztásokat, tudásalapú szabályokkal értékelte őket, majd korrelálta az eredményeket~\cite{yu2005trinetr}; Taha és szerzőtársai szintén agentekre bontották a riasztások korrelációját~\cite{taha2010}. Ezekben a rendszerekben az agent szerepe a munka megosztása és a modularitás volt; a döntések szabályokon alapultak, ezért kiszámíthatók, de a szabályok karbantartása nagy munkát igényelt.

\subsubsection{A nagy nyelvi modellek megjelenése a SOC-ban}
A nagy nyelvi modellek (\ac{LLM}) új lehetőséget hoztak: természetes nyelvű szöveget értenek és állítanak elő, így képesek a riasztások tartalmát összefoglalni, a vizsgálat következő lépését javasolni, vagy a szervezet dokumentációjából releváns részeket előkeresni. Ismail és szerzőtársai például a Wazuh eseményeihez olyan segédet készítettek, amely a választ egy keresésen alapuló szövegelőállítással (\ac{RAG}) támasztja alá: az incidenskezelési tudásból, a MITRE ATT\&CK tudásbázisból és a NIST kiberbiztonsági keretrendszeréből keres elő összefüggéseket~\cite{ismail2025copilot}. Hmimou és szerzőtársai három szakosodott agentet -- levélellenőrző, naplóelemző és IP-cím-vizsgáló -- kötöttek össze egy központi összetevővel, amely az eredményeiket nyelvi modell segítségével korrelálja~\cite{hmimou2025}. Chhetri és szerzőtársai hangsúlyozzák, hogy ezek a rendszerek az elemzőt nem kiváltják, hanem vele együttműködve kell működniük~\cite{chhetri2024}.

\subsubsection{Nyelvi modellekre épülő agentrendszerek a riasztások kezelésében}
A legújabb munkák több, nyelvi modellel működő agentet kapcsolnak össze a riasztások vizsgálatára. A CORTEX-ben egy-egy agent a viselkedést elemzi, a bizonyítékokat gyűjti és a döntést hozza meg; az adatokhoz szabványos eszközfelületen (\ac{MCP}) férnek hozzá~\cite{wei2025cortex}. A szerzők mérése szerint a többágenses változat pontosabb volt az egyetlen modellnél: a beavatkozást igénylő esetek felismerésének F1-értéke 0,66-ról 0,78-ra nőtt, a téves riasztások aránya pedig 24,9\%-ról 14,2\%-ra csökkent. Ennek ára azonban számottevő: riasztásonként mintegy 23600 szövegegységet (tokent) dolgozott fel, nagyjából 5,7-szer többet, és egy riasztás vizsgálata átlagosan 152 másodpercig tartott, ami körülbelül 3,4-szeres lassulás.

A LanG egységes biztonsági műveleti platformot ír le, amelyben a korrelációt gráfalapú közösségkereső algoritmus és valószínűségi pontozás végzi, a nyelvi modell pedig csak hipotéziseket fogalmaz meg; a bizonytalan válaszokat egy küszöb kiszűri, a lényeges döntésekhez pedig emberi jóváhagyás kell~\cite{abdennebi2026lang}. Az AgentSOC hasonló elven a nyelvi modell javaslatait egy kiszámítható, gráfalapú ellenőrzésen engedi át, és a gépnyilvántartás adatait is felhasználja; kiértékelése azonban egy kis méretű, részben mesterséges környezetben, előre felcímkézett összehasonlító adathalmaz nélkül történt, és a futásidő nagyjából 95\%-át a nyelvi modell tette ki~\cite{roy2026agentsoc}.

\subsubsection{A nyelvi modellek korlátai és a vegyes felépítés}
A nyelvi modellek használatának több, jól dokumentált korlátja van. Kimenetük nem determinisztikus, így ugyanarra a riasztásra két futás eltérő választ adhat; ez a megismételhető mérést és az utólagos elszámoltathatóságot is megnehezíti. Hihető, de hamis állításokat is előállíthatnak, a feldolgozható szöveg hossza korlátos, és a futtatásuk -- különösen nagy riasztásmennyiségnél -- lassú és költséges. A felhőben futó modellek használata ráadásul azt jelenti, hogy érzékeny naplóadatok hagyják el a szervezetet.

Liu és Wang vizsgálata szerint a helyben futtatható, kisebb modellek önállóan nem képesek megbízhatóan végigvezetni egy riasztás vizsgálatát: rossz lépést választanak, vagy nem tudják, mikor kell megállniuk~\cite{liu2026hesp}. Ha viszont a lépések sorrendjét és a leállás feltételét a modellen kívüli, kiszámítható vezérlő határozza meg, egy hétmilliárd paraméteres modell ellenőrzött befejezési aránya 0,125-ről 1,000-ra nőtt. A SENTINEL-RL hasonló következtetésre jut: a hálózati összefüggések elemzését és a döntést egy erre tanított, nem nyelvi modell végzi, a nyelvi modell csak az eredmény szöveges magyarázatát írja meg, amelyet egy további ellenőrző lépés hagy jóvá~\cite{vallabhaneni2026sentinel}.

Ezekből a munkákból egy közös irány rajzolódik ki: a nyelvi modell ott hasznos, ahol szöveget kell értelmezni vagy előállítani, a pontozást és a lépések sorrendjét azonban célszerű kiszámítható programra bízni. A jelenlegi nyelvimodell-alapú rendszerek közös gyengesége ugyanakkor az értékelés: a legtöbb saját, nem nyilvános vagy mesterségesen előállított adaton mér, ezért eredményeik nem vethetők össze egymással, és nem derül ki, mennyit adnak hozzá egy erős, nyelvi modell nélküli rangsoroláshoz.

\begin{table}[H]
    \centering
    \caption{Agent-alapú és nyelvimodell-alapú riasztáskezelő rendszerek (Forrás: saját összeállítás)}
    \label{tab:agentek}
    \small
    \begin{tabular}{|L{2.3cm}|L{3.2cm}|L{3.0cm}|L{4.7cm}|}
        \hline
        \textbf{Rendszer} & \textbf{A nyelvi modell szerepe} & \textbf{Kiértékelés} & \textbf{Gyengeség} \\ \hline
        TRINETR \cite{yu2005trinetr} & nincs (szabályok) & saját környezet & Kézi szabálykarbantartás \\ \hline
        CORTEX \cite{wei2025cortex} & minden lépés & saját adat & Kb. 5,7-szeres szövegmennyiség, 3,4-szeres lassulás \\ \hline
        LanG \cite{abdennebi2026lang} & hipotézis, küszöbbel szűrve & saját adat & Összetett platform; nincs nyilvános összehasonlító mérés \\ \hline
        AgentSOC \cite{roy2026agentsoc} & hipotézis, gráffal ellenőrizve & kis, részben mesterséges környezet & Nincs címkézett összehasonlító adathalmaz \\ \hline
        HESP \cite{liu2026hesp} & vizsgálati lépések, külső vezérlővel & saját feladatsor & A rangsor minőségét nem méri \\ \hline
        SENTINEL-RL \cite{vallabhaneni2026sentinel} & csak magyarázat, ellenőrzéssel & mesterséges hálózat & Tanított döntési modellt igényel \\ \hline
    \end{tabular}
\end{table}

\subsection{A kiértékelés gyakorlata és buktatói}\label{sec:ertekeles-irodalom}
A módszerek összehasonlítását nehezíti, hogy kiértékelésük nagyon eltérő módon történik. Arp és szerzőtársai a gépi tanulást alkalmazó biztonsági kutatások gyakori hibáit gyűjtötték össze~\cite{arp2022}. Ezek közül a riasztáskezelés szempontjából a legfontosabbak: a valóságot rosszul tükröző mintavétel; az, ha a módszer beállításai a tesztadatból származó tudást is felhasználnak; a túl gyenge összehasonlítási alap; a feladathoz nem illő mérőszám; a valódi támadások ritkaságának figyelmen kívül hagyása; valamint az, ha a módszert csak laboratóriumi körülmények között vizsgálják. Ndichu és szerzőtársai áttekintése szerint a riasztási fáradtság csökkentésére javasolt módszerek közül 87 tanulmányból csak 9 alapult éles SOC-adatokon~\cite{ndichu2026}.

A megismételhető összehasonlításhoz nyilvános, felcímkézett adathalmaz kell. Az AIT-ADS egy több lépésből álló támadást nyolc, eltérő felépítésű környezetben játszik le, és a hálózati és gépi érzékelők riasztásait a támadáshoz tartozás jelölésével együtt teszi közzé~\cite{landauer2024aitads}. Uetz és szerzőtársai ugyanezen az adathalmazon mérték a kockázatalapú pontozást~\cite{uetz2026rba}, így a dolgozat eredményei közvetlenül összevethetők az övékkel. A dolgozat értékelési terve ezekre a tanulságokra épül: a szakértői súlyokat a mérés előtt rögzítem, az adatból tanult súlyokat egy forgatókönyv kihagyásával tanítom és azon mérem, és az összehasonlítási alap a legerősebb nyilvánosan mért módszer.

\begin{table}[H]
    \centering
    \caption{A meglévő megoldások hiányosságai és a tervezett rendszer válaszai (Forrás: saját összeállítás)}
    \label{tab:resek}
    \small
    \begin{tabular}{|L{4.4cm}|L{2.6cm}|L{6.6cm}|}
        \hline
        \textbf{Hiányosság} & \textbf{Érintett módszerek} & \textbf{A tervezett rendszer válasza} \\ \hline
        A legerősebb nyilvános, tanítás nélküli módszer nem használ környezeti információt & \cite{uetz2026rba} & A kockázatalapú pontszámot a gép fontosságával, a sérülékenységekkel és a megfigyeltséggel bővíti, lépésenként mérve a hozzáadott értéket \\ \hline
        A környezeti módszerek szűrnek, hiányos nyilvántartásnál valódi támadást is eltávolítanak & \cite{njogu2012,sommestad2015} & A környezet csak súlyként módosítja a sorrendet, riasztást nem töröl; az átcsúszó támadások arányát legfeljebb 2\%-ban korlátozza \\ \hline
        A szakértői súlyok hatását ritkán vizsgálják & \cite{anuar2012,alsubhi2011fuzmet} & Páros összehasonlítással (\ac{AHP}) rögzített súlyok, érzékenységvizsgálat és adatból tanult súlyokkal való összevetés \\ \hline
        A tanuló módszerek címkét vagy részletes gépi naplót igényelnek & \cite{wang2024alertpro,hassan2019} & Tanítóadat nélkül működik; tanult súlyokat csak összehasonlításként használ \\ \hline
        A nyelvimodell-alapú agentek lassúak és nem megismételhetők & \cite{wei2025cortex,liu2026hesp} & A pontszámot kiszámítható program adja, a lépések sorrendjét vezérlő tartja; a nyelvi modell csak magyarázatot ír, helyben futtatva, ellenőrzéssel \\ \hline
        Az eredmények nem nyilvános vagy mesterséges adatokon születnek & \cite{ndichu2026,roy2026agentsoc} & Nyilvános, felcímkézett adathalmazon (AIT-ADS) mér, a meglévő mérésekkel közvetlenül összevethetően \\ \hline
        A csoportosítás és a rangsorolás külön-külön készül & \cite{landauer2022,uetz2026rba} & Az összetartozó riasztásokat csoportosítja, és a csoportokat rangsorolja \\ \hline
    \end{tabular}
\end{table}